\section{Introducción}
\label{ch:introduccion}

En los últimos años ha surgido la necesidad de digitalizar la mayoría de los procesos, por lo que los sistemas de empresas gubernamentales atraviesan un proceso de digitalización, en el cual hay muchos módulos que requieren complementar su sistema para garantizar la eficiencia y agilidad en los procesos administrativos. El Poder Judicial del Estado de Tamaulipas (PJETAM) tiene una gran cantidad de archivos que requieren de urgente organización para asegurar un acceso rápido y un almacenamiento adecuado. El proyecto abordado en esta tesina surge como una respuesta tecnológica para modernizar la forma en que manejan y acceden a los archivos en las distintas áreas de la organización del PJETAM.

El principal propósito de este proyecto es documentar el desarrollo e implementación de una aplicación web capaz de manejar la carga, conversión, almacenamiento y gestión de documentos, a través del uso de tecnologías modernas para crear una interfaz completa, además de almacenamiento en la nube. El proyecto busca crear una aplicación utilizando la arquitectura MVC  (Modelo-Vista-Controlador) para lograr una aplicación segura, escalable y eficiente.

Para lograr el objetivo, se busca realizar una aplicación que permita a las diferentes áreas del PJETAM cargar, convertir, almacenar y gestionar documentos e imágenes, generando códigos QR vinculados a cada archivo para facilitar su referenciación, e integrándose con los sistemas de escritorio existentes mediante autenticación por parámetro cifrado en la URL. Para realizar este proyecto se tienen objetivos específicos, que abarcan desde el análisis de requerimientos funcionales y el diseño de la arquitectura tecnológica, hasta la implementación de módulos de conversión automática de archivos ODT, DOC, DOCX a PDF, implementación de mecanismos de seguridad y la documentación del proyecto.

El propósito de este trabajo es optimizar la eficacia en la gestión de documentos. Ya que actualmente no existe un sistema que haga todo lo mencionado anteriormente, por lo que es fundamental hacer un sistema que cumpla estos requerimientos. La adopción de un sistema con conversión automática a PDF e integración utilizando códigos QR ayuda enormemente a reducir la cantidad de tiempo dedicada a la búsqueda de información, y facilita la navegación entre documentos. Además, simplifica la transición de formatos de documento de papel al digital.

El alcance de este proyecto se centra en el desarrollo, implementación y pruebas de una aplicación web para la gestión de documentos orientada únicamente a los trabajadores del Poder Judicial del Estado de Tamaulipas (PJETAM). Las limitaciones del proyecto se centran en el acceso al público externo, ya que estará estrictamente limitado a la visualización de aquellos archivos que sean marcados como públicos o puestos a disposición del público. En la parte del desarrollo una limitación es que los procesos de seguridad funcionarán temporalmente usando la autenticación URL cifrada procedente de otros módulos actuales.

Para detallar la elaboración de este proyecto, este documento se estructura en las siguientes secciones: Introducción, que aborda el planteamiento del problema y los objetivos del proyecto, después se presenta la sección de Metodología, que describe la metodología de desarrollo y diseño de la arquitectura del sistema; y finalmente, se presentan los resultados obtenidos, las pruebas de integración y las conclusiones del proyecto.

\subsection{Antecedentes}
\label{sec:antecedentes}
En esta sección, se presenta una revisión concisa y pertinente de investigaciones previas, literatura existente y el estado del arte relacionado con el tema de la tesis. El objetivo es situar el trabajo actual dentro del conocimiento acumulado, identificar vacíos en la investigación y demostrar la originalidad y la necesidad del estudio propuesto. Se deben citar adecuadamente las fuentes, mostrando cómo el trabajo se construye sobre cimientos previos o cómo aborda aspectos no resueltos.

\subsubsection{Antecedentes de la Empresa}
\label{subsec:antecedentes_empresa}
El Supremo Tribunal de Justicia del Estado de Tamaulipas, ubicado en Boulevard Praxedis Balboa número 2207, entre López Velarde y Díaz Mirón, Colonia Miguel Hidalgo, Código Postal 87090, en Ciudad Victoria, Tamaulipas, México. Formalizó su inicio el 9 de julio de 1824 con la expedición del primer decreto de la Legislatura Constituyente local.

La misión de la empresa es Impartir y administrar justicia de manera eficiente y expedita, solucionando conflictos a través de los órganos jurisdiccionales y los medios alternos, contribuyendo a garantizar el estado de derecho y la paz social en Tamaulipas.

La visión de la empresa es Transcender como Institución de excelencia, mediante la profesionalización e institucionalización de sus integrantes, en un ambiente de humanismo y justicia institucional, haciendo uso de la innovación tecnológica y optimización de los recursos, con responsabilidad social y pleno respeto a la ley, los derechos humanos y la igualdad de género.

Estoy realizando las estadías TSU en el departamento de informática, realizando el proyecto Sistema Web de Gestión Documental y Generación de Códigos QR para el PJETAM, el cual consiste en un sistema diseñado para el personal de la empresa, para una mejor organizacion y almacenamiento de los documentos. 

El proyecto presentado en esta tesina es un proyecto desde cero, por lo cual antes de este proyecto la busqueda de documentos y almacenamiento ordenado de los distintos archivos se debía realizar en multiples computadoras y buscar el documento en distintas paginas de la empresa.

\subsubsection{Antecedentes Teoricos}
\label{subsec:antecedentes_teoricos}

\subsection{Definición del Problema}
\label{sec:problema}

El manejo de la información y archivos dentro de las instituciones públicas es un elemento esencial para el cumplimiento eficiente de sus tareas. Sin embargo, el Poder Judicial del Estado de Tamaulipas (PJETAM) actualmente presenta una brecha tecnológica en la administración de sus archivos recurrentes. El problema central consiste en la desorganización de los documentos institucionales y la dificultad que presentan los sistemas actuales para proporcionar una navegación fácil y rápida al consultar expedientes o documentos. Tales factores generan un entorno insatisfactorio que tiene consecuencias directas sobre el desempeño de las áreas jurídicas. 

Del problema mencionado se derivan distintas complicaciones que no han sido abordadas individualmente:

\begin{itemize}
    \item \textbf{Falta de estandarización de formatos:} Los documentos son generados en múltiples formatos (ODT, DOC, DOCX), lo que provoca incompatibilidades y problemas de formato al compartirlos entre distintas áreas o al intentar visualizarlos rápidamente.
    \item \textbf{Desconexión entre el documento físico y el digital:} Los procesos judiciales a menudo requieren referenciar documentos digitales en expedientes físicos. Actualmente, buscar el formato digital de un documento impreso es un proceso lento y manual, propenso a errores humanos.
    \item \textbf{Riesgo de pérdida o duplicidad de información:} La ausencia de un sistema centralizado de almacenamiento fomenta que los archivos se guarden de manera local o fragmentada, dificultando su organización, consulta y eventual eliminación o depuración.
\end{itemize}

La falta de una herramienta que convierta y organice estos archivos es igual a horas de trabajo invertidas en tareas puramente administrativas y de búsqueda, en lugar de otras labores de mayor importancia.

Ante este contexto, se plantea la siguiente pregunta de investigación para guiar el desarrollo de este proyecto:
¿De qué manera el desarrollo e implementación de un Sistema Web de gestión documental, con capacidades de conversión automática de formatos y generación de códigos QR, mejorará la organización, el acceso y la interoperabilidad de los documentos en las diferentes áreas del PJETAM?

\subsection{Objetivos}
\label{sec:objetivos}

El objetivo general del proyecto presenta la meta final del proyecto:
\subsubsection{Objetivo General}
\label{subsec:objetivo_general}

Desarrollar una aplicación web que permita a las diferentes áreas del Poder Judicial del Estado de Tamaulipas cargar, convertir, almacenar y gestionar documentos e imágenes, generando códigos QR vinculados a cada archivo para facilitar su referenciación y consulta desde otros documentos institucionales. La aplicación deberá integrarse con los sistemas de gestión de escritorio existentes mediante autenticación por parámetro cifrado en la URL.

\subsubsection{Objetivos Específicos}
\label{subsec:objetivos_especificos}

Los objetivos especificos de este proyecto son los siguientes:
\begin{itemize}
    \item Analizar los requerimientos funcionales de las áreas usuarias del PJETAM, identificando los tipos de documentos, flujos de trabajo y necesidades de integración con los sistemas de escritorio existentes.
    \item Diseñar la arquitectura de la aplicación web, definiendo las tecnologías de frontend, backend y almacenamiento que garanticen seguridad, rendimiento y escalabilidad.
    \item Implementar el módulo de carga y conversión de archivos, permitiendo recibir documentos en formatos ODT, DOC y DOCX para su conversión automática a PDF, así como imágenes y URLs.
    \item Desarrollar el módulo de almacenamiento y gestión documental, que permita guardar, consultar y eliminar los archivos registrados en el sistema.
    \item Implementar la generación automática de códigos QR vinculados a cada documento o archivo almacenado, permitiendo su descarga e inserción en documentos físicos o digitales de las diferentes áreas jurídicas.
    \item Desarrollar el mecanismo de autenticación mediante parámetro cifrado en la URL, que permita la apertura segura de la aplicación web desde los distintos sistemas de gestión de escritorio del PJETAM.
    \item Realizar pruebas de integración con los sistemas de escritorio existentes, verificando el correcto funcionamiento del acceso autenticado y la interoperabilidad entre plataformas.
    \item Elaborar la documentación técnica y manuales de usuario para la operación, administración y mantenimiento de la aplicación.
\end{itemize}

\subsection{Justificación}
\label{sec:justificacion}
La administración del PJETAM requiere procesos administrativos rápidos, claros y altamente organizados, ya que es una institución que maneja procedimientos legales y archivos importantes. El manejo de archivos es una funcion básica y fundamental del PJETAM. La realización de este proyecto se justifica por la necesidad de resolver la desorganización de los archivos institucionales y la falta de agilidad en su consulta, factores que actualmente aumentan los tiempos de respuesta de la organización y complican el flujo de trabajo de los empleados.

El desarrollo del Sistema Web de Gestión Documental y Generación de Códigos QR aporta un beneficio a la institución, ya que uno de los problemas de la organización es la diversidad de formatos en archivos (ODT, DOC, DOCX) y la división de la información que generan un conflicto al momento de visualizar y compartir documentos. Al implementar un sistema de conversión automática a PDF el proyecto garantiza la estandarización de los archivos, asegurando que cualquier documento se visualice de igual manera sin importar el dispositivo o el área que lo consulte, asimismo imparte un estándar al almacenar todos los archivos en un mismo formato. Además, la generación de códigos QR resuelve un problema práctico el cual es el referenciar archivos en otro archivo, haciendo más fácil la navegación entre documentos y facilitando la lectura y comprensión del lector.
\subsection{Alcances y Limitaciones}
\label{sec:alcances_limitaciones}

\subsubsection{Alcances}
\label{subsec:alcances}
El alcance de este proyecto incluye el desarrollo, implementación y pruebas de una aplicación web de gestión de archivos destinada al personal del Poder Judicial del Estado de Tamaulipas (PJETAM). La aplicación se centrará en crear un módulo para cargar archivos y gestionar su almacenamiento, lo que implicará usar Amazon Web Services (AWS) S3 y bases de datos SQL Server. 
La aplicación ofrecerá la conversión automatizada de  distintos tipos de archivos (como ODT, DOC, DOCX) a formato PDF, así como imágenes y URLs, una vez guardados los archivos con un mismo formato se podran gestionar en el modulo de archivos, en el cual permitira obtener codigos QR vinculado a cada archivo asi como la opcion de descargarlo o elminarlo. Uno de los aspectos clave es la creación automática de códigos QR para cada documento almacenado, lo que permitirá su inserción y consulta de forma rápida tanto desde expedientes físicos como digitales. Además, la plataforma web estará diseñada para integrarse con los sistemas de gestión de escritorio existentes, asegurando la interoperabilidad mediante un mecanismo de autenticación basado en parámetros cifrados en la URL.

\subsubsection{Limitaciones}
\label{subsec:limitaciones}
Las limitaciones de este proyecto son la falta de implementación de un sistema de roles individuales, ya que el sistema sólo registrará el área a la que pertenece cada movimiento realizado en la aplicación, ya que así será organizado el almacenamiento de archivos. Otra limitación es que el sistema no contará con un módulo de edición de texto en documentos. 
En el ámbito de desarrollo, una limitación es la creación de módulos de seguridad complejos adicionales a la URL cifrada; de igual forma, el acceso al público externo estará limitado a la visualización de aquellos archivos que sean puestos a disposición de la ciudadanía.

\subsection{Ejemplos de Tablas e Inclusión de Imágenes}
\label{sec:ejemplos_elementos}

Para ilustrar el uso de elementos visuales en su tesis, a continuación se presentan ejemplos básicos de cómo incluir tablas e imágenes en LaTeX.

\subsubsection{Inclusión de Tablas}

Las tablas son herramientas fundamentales para presentar datos de manera organizada y comprensible.
\begin{table}[H]
    \centering
    \caption{Clasificación de Sensores Comunes en Ingeniería.}
    \label{tab:sensores_comunes}
    \begin{tabular}{|l|c|c|}
        \hline
        \textbf{Tipo de Sensor} & \textbf{Magnitud Medida} & \textbf{Principio de Funcionamiento} \\
        \hline
        Termopar & Temperatura & Efecto Seebeck \\
        Extensómetro & Deformación & Resistencia eléctrica \\
        Acelerómetro & Aceleración & Efecto piezoeléctrico \\
        Encoder rotatorio & Posición angular & Óptico/Magnético \\
        \hline
    \end{tabular}
    \caption*{Nota: Esta tabla es un ejemplo para demostrar la estructura básica de una tabla en LaTeX.}
\end{table}
Como se puede observar en la Tabla \ref{tab:sensores_comunes}, la información se presenta de forma clara y estructurada. Es importante añadir un \verb|`\caption`|  para describir la tabla y un  \verb|`\label`|  para referenciarla en el texto.

\subsubsection{Inclusión de Imágenes}

Las imágenes, gráficos o diagramas son cruciales para complementar la explicación textual y facilitar la comprensión de conceptos complejos o resultados visuales.
\begin{figure}[H]
    \centering
    \includegraphics[width=0.7\textwidth]{img/diagrama_proceso}
    \caption{Diagrama de un proceso de control automático.}
    \label{fig:diagrama_control}
\end{figure}
% La Figura \ref{fig:diagrama_control} ilustra un ejemplo de un diagrama de un proceso de control. Para incluir una imagen, se utiliza el entorno `figure` y el comando `\includegraphics`. Es fundamental especificar la ruta de la imagen (en este caso, `images/diagrama_proceso.png`) y un `width` para controlar su tamaño. Al igual que con las tablas, un `\caption` y un `\label` son indispensables. Es crucial que cada imagen tenga una calidad adecuada y que sea relevante para el texto.

\begin{figure}[H]
    \centering
    \begin{subfigure}[b]{0.45\textwidth}
        \includegraphics[width=\textwidth]{img/grafico_datos1}
        \caption{Datos de Temperatura.}
        \label{fig:subfig_temp}
    \end{subfigure}
    \hfill
    \begin{subfigure}[b]{0.45\textwidth}
        \includegraphics[width=\textwidth]{img/grafico_datos2}
        \caption{Datos de Presión.}
        \label{fig:subfig_pres}
    \end{subfigure}
    \caption{Comparación de datos operativos de un sistema (Temperatura vs. Presión).}
    \label{fig:comparacion_datos}
\end{figure}
Incluso es posible combinar varias imágenes en una sola figura, como se muestra en la Figura \ref{fig:comparacion_datos}, utilizando el paquete `subcaption`.

Asegúrense de que todas las figuras y tablas estén debidamente referenciadas en el texto y que su contenido sea autoexplicativo, utilizando los pies de figura y títulos de tabla de manera efectiva.